%%%%%%%%%%%%%%%%%%%%%%%%%%%%%%%%%%%%%%%%%%
% Mathematics Final Year Research Projects
% LaTeX Template
% Version 1.0 (31/01/24)
%
% This template has been adapted from: https://www.overleaf.com/latex/templates/imperial-college-report-template/wncnzptkhnbc
% Students should feel free to adapt this template to their needs.
%%%%%%%%%%%%%%%%%%%%%%%%%%%%%%%%%%%%%%%%%%
%----------------------------------------------------------------------------------------
% PACKAGES AND OTHER DOCUMENT CONFIGURATIONS
%----------------------------------------------------------------------------------------
\documentclass[a4paper,11pt, twoside]{report}

%% Language and font encodings
\usepackage[english]{babel}
\usepackage[utf8]{inputenc}
\usepackage[T1]{fontenc}

%% Sets page size and margins
\usepackage[a4paper,top=1in,bottom=1in,left=1in,right=1in,marginparwidth=1.75cm]{geometry}

%% Useful packages
\usepackage{afterpage}

\usepackage{amsmath}
\usepackage{amsthm}
\usepackage{amssymb}
\usepackage{csquotes}
\usepackage{enumitem}
\usepackage{graphicx}
\usepackage{lipsum}
\usepackage{booktabs}
\usepackage{multirow}
\usepackage{longtable}
\usepackage{bbm}
% Listings (for displaying code):
\usepackage{listings}
\lstset{
    frame = single, 
    framexleftmargin=15pt
}

% Center figure captions:
\usepackage{caption}
\captionsetup[figure]{labelfont={bf},name={Figure},labelsep=quad}
\captionsetup[table]{labelfont={bf},name={Table},labelsep=quad}


% Theorems
\newtheorem{definition}{Definition}
\newtheorem{theorem}{Theorem}
\newtheorem{proposition}{Proposition}
\newtheorem{corollary}{Corollary}
\newtheorem{lemma}{Lemma}

% ----------- Algorithm2e setup
\usepackage[ruled,vlined]{algorithm2e}
\makeatletter
\renewcommand{\SetKwInOut}[2]{%
  \sbox\algocf@inoutbox{\KwSty{#2}\algocf@typo:}%
  \expandafter\ifx\csname InOutSizeDefined\endcsname\relax% if first time used
    \newcommand\InOutSizeDefined{}\setlength{\inoutsize}{\wd\algocf@inoutbox}%
    \sbox\algocf@inoutbox{\parbox[t]{\inoutsize}{\KwSty{#2}\algocf@typo:\hfill}~}\setlength{\inoutindent}{\wd\algocf@inoutbox}%
  \else% else keep the larger dimension
    \ifdim\wd\algocf@inoutbox>\inoutsize%
    \setlength{\inoutsize}{\wd\algocf@inoutbox}%
    \sbox\algocf@inoutbox{\parbox[t]{\inoutsize}{\KwSty{#2}\algocf@typo:\hfill}~}\setlength{\inoutindent}{\wd\algocf@inoutbox}%
    \fi%
  \fi% the dimension of the box is now defined.
  \algocf@newcommand{#1}[1]{%
    \ifthenelse{\boolean{algocf@inoutnumbered}}{\relax}{\everypar={\relax}}%
%     {\let\\\algocf@newinout\hangindent=\wd\algocf@inoutbox\hangafter=1\parbox[t]{\inoutsize}{\KwSty{#2}\algocf@typo\hfill:}~##1\par}%
    {\let\\\algocf@newinout\hangindent=\inoutindent\hangafter=1\parbox[t]{\inoutsize}{\KwSty{#2}\algocf@typo:\hfill}~##1\par}%
    \algocf@linesnumbered% reset the numbering of the lines
  }}%
\makeatother
% --------- end algorithm2e setup

% \bm allows typing bold math:
\usepackage{bm}
\usepackage[normalem]{ulem}

\usepackage[colorinlistoftodos]{todonotes}
\usepackage[colorlinks=true, allcolors=blue]{hyperref}

\renewcommand*{\rmdefault}{bch}
\renewcommand*{\ttdefault}{lmtt}
\newcommand{\citationneeded}{\textcolor{red}{[citation-needed]}}

\DeclareMathOperator*{\argmin}{\arg\!\min}
\DeclareMathOperator*{\argmax}{\arg\!\max}

% Add bigger skip between paragraphs, makes reading easier:
\setlength{\parskip}{0.5em}

% Bibliography
\usepackage[numbers, comma, square, sort&compress]{natbib}
\bibliographystyle{abbrvunsrtnat.bst}

%----------------------------------------------------------------------------------------
% END OF DOCUMENT CONFIGURATION
%----------------------------------------------------------------------------------------

%----------------------------------------------------------------------------------------
% IMPORTANT INFORMATION TO MODIFY FOR THE TITLE PAGE
%----------------------------------------------------------------------------------------
\newcommand{\reporttitle}{Bayesian Phylolinguistic Inference under Cognate Dependence: A Study of Model Misspecification} % Title of your research project
\newcommand{\reportauthor}{Yong Sheng James Tay} % First Name and Last Name
\newcommand{\supervisor}{Dr Robin Ryder} % First Name and Last Name of your supervisor(s)
\newcommand{\degreetype}{MSci in Mathematics} % MSci in Mathematics, BSc in Mathematics, BSc in Mathematics with Statistics, ...
%----------------------------------------------------------------------------------------

%----------------------------------------------------------------------------------------
% To compile this file, use the following sequence: 
%	latex main.tex
%	bibtex main.tex
%	latex main.tex
%	latex main.tex
%----------------------------------------------------------------------------------------

%----------------------------------------------------------------------------------------
% START OF DOCUMENT
%----------------------------------------------------------------------------------------
\begin{document}

% Title page
\input{title/title.tex}

% Abstract
\begin{abstract}
In historical linguistics, Bayesian phylogenetic methods have quickly become the principal framework for reconstructing the evolution of languages from cognate data. One key assumption of such methods is conditional site independence: that cognacy classes evolve independently from each other. However, given the relative complexities of language evolution, linguistic and cultural datasets often exhibit dependencies between sites. This causes model misspecification: a mismatch in assumptions made by the model and the true characteristics of the data, the consequences for which remain poorly understood.

This thesis investigates the impact of violating the conditional site independence assumption on Bayesian phylogenetic inference, applied in linguistic contexts. A simulation study was conducted to answer two questions: (1) the extent by which unknown dependencies amongst sites of the data had negative effects on tree reconstruction, and (2) whether including known correlated sites (and thus causing model misspecification) affected tree reconstruction. These questions were posed to provide an understanding of the implications caused by such dependency structures and a framework for decision making for retaining dependent data.

The results show that violating conditional site independence indeed reduces tree reconstruction quality relative to fully independent data. The depth of dependency, however, made little contribution to affecting the tree reconstruction quality. Instead, the greatest factor of reconstruction quality is the total \textit{effective} information contained within the dataset. In low-information scenarios, unknown dependencies cause larger deterioration while additional dependent data generally improves reconstruction despite violating model assumptions. On the other hand, in high-information scenarios, the drop in reconstruction quality due to unknown dependencies is minor, while marginal benefits of additional dependent sites diminish and thus may not justify their inclusion.

These findings identify effective information as the underlying framework for understanding the effects of model misspecification in Bayesian phylolinguistics. They suggest that the implications of violating conditional site independence in such methods are a nuanced matter, where correlated sites present tradeoffs between introducing bias and additional effective information. Depending on the context at hand, these factors then determine whether their inclusion proves beneficial for phylogenetic reconstruction.

\end{abstract}
 
% Acknowledgements
\clearpage
\section*{Acknowledgments}
This thesis would not have been possible without the help and support of many.

Above all, I thank the Lord for His goodness in my life, not least in the completion of this thesis. Mathematics and Linguistics are but two of the many \textit{good and perfect} gifts God has given us; our study of their complexities and intricacies serves only to reflect, in part, the order and creativity of His creation.

I would like to express my sincere gratitude to my supervisor, Dr Robin Ryder, for his guidance and support throughout the course of this thesis. It has been a true privilege to have worked under his tutelage, not only through his insightful guidance and advice, which have been pivotal at key moments, but also because of his enthusiasm for historical linguistics. Our conversations have been invaluable in shaping not only this project but also my understanding of the field itself. 

I thank Guillaume Jacques and Chris Buckley for their insights on the results of the thesis and their application in linguistic data. I am also grateful to Imperial College London's High Performance Computing service for providing the computational resources required for the simulations performed in this thesis.

My family have been a source for great comfort. I thank my parents for their unconditional love and support throughout the course of my education. I may never be able to comprehend the extent of your sacrifices and trust you gave me to pursue my academic passions freely. I thank my brother for being an inspiration and, at times, a grounding presence, as we spur each other on in our pursuits.

Finally, my undergraduate studies have been some of the best and most formative years of my life, for which I have my friends to thank. In particular, I thank DB, AJ, CM, YH, LN, GC, DK and EM for being constant presences in my time at university. Your love and encouragement have had real impact on my life which I will always cherish.

\clearpage

% Acknowledgements
\clearpage
\section*{Plagiarism statement}
The work contained in this thesis is my own work unless otherwise stated. \\

\vspace{1em}
\noindent \textit{Signature:} \reportauthor \\
\textit{Date:} \today
\clearpage

% Table of contents
\tableofcontents

% Main sections of the report
% ===========================
% Uncomment or add folders to add your own chapters and input files.

\chapter{Introduction}

The introduction should attempt to set your work in the context of other work done in the field. It
should demonstrate that you are aware of what you are doing, and how it relates to other work
(with references). It should also provide an overview of the contents of the project. You should
highlight your individual contributions and any novel result: which of the calculations, theorems,
examples, proofs, conjectures, codes etc. are your own?

Language diversification and phylogenetic inference
\section{History of Language Diversification Methods}
\section{Bayesian Phylogenetic Methods}
    Bayesian phylogenetic methods
\subsection{Site Independence}
    The independence assumption
    Why independence may fail in linguistic data
\section{Previous Work}
\subsubsection{In Comparative Biology}
            Maddison's Tail Colour Problem
            Longobardi Correlated Issues
\subsubsection{In Phylogenetic Inference}
        In phylogenetic inference
            see examples in readings file
\subsubsection{In Linguistic Data}
        Dependence in linguistic data
            Kra-dai data
\subsection{Outstanding gap in Literature}
        Outstanding gap in literature
\section{Research Questions and Contributions}
    Research questions and contributions (Novel work)
        (1) how does an increased proportion of fixed amounts of data being dependent (on the remaining independent data) affect 
        reconstruction? if we didnt know how much data was actually dependent and assumed independence, how badly does this 
        affect reconstruction? 
        (2) how does adding dependent data to a set amount of independent data affect reconstruction? is it worth adding more 
        data even though the data becomes misspecified by the model? 

\section{Thesis Structure}


\section{ (TBD) helpful things (subequations, algorithms)}
To write equations over multiple lines (like systems of equation or equations too long to fit the page), one can use the \textit{align} environment coupled to the \textit{subequations} environment like so
\begin{subequations}
\begin{align}
    \mathbb{P}(0 \leadsto -a) &= D\int_0^\infty dt\, k_r e^{-k_r t}\int_0^t d\tau\,\partial_x u \big|_{x=-a}, \\
    \mathbb{P}(0 \leadsto b) &= -D\int_0^\infty dt\, k_r e^{-k_r t}\int_0^t d\tau\,\partial_x u \big|_{x=b}.
\end{align}
\end{subequations}

\begin{algorithm}[]
\SetAlgoLined
\SetKwInOut{KwInput}{Input}
\SetKwInOut{KwOutput}{Output}
\SetKwInOut{KwPre}{Pre}
\SetKw{Return}{return}
\SetKwProg{Fn}{Function}{}{end}
\LinesNumbered
\KwInput{$\textbf{m}$, such that $\mathbf{m}_i$ is the position of the $i$'th monitor\newline
$\textbf{l}$, such that $\mathbf{l}_i$ is the position of the $i$'th landmark\newline
$\mathbf{p}^m$, such that $\mathbf{p}^m_i$ is the ping latency from monitor $i$ to the target\newline
$\mathbf{p}^l$, such that $\mathbf{p}^l_i$ is the set of ping latencies to landmark $i$}

\BlankLine
\KwPre{Compute $\hat{p}_m(d \mid l)$, an estimator giving the likelihood of the target being distance $d$ away from the monitor, given that the monitor records a latency of $l$ to that target. Implemented by training a KDE using $\mathbf{p}^l$.\newline
Compute $\hat{p}_l(d \mid v)$, an estimator giving the likelihood of the target being distance $d$ away from the landmark, given a Canberra distance of $v$ between the target and the landmark, using training targets.
}
\BlankLine
\KwOutput{Most likely location of the target}
\BlankLine

\Fn{Likelihood($x$, $\mathbf{v}$)} {
MonitorScore $\gets \sum_{i=0}^{|\mathbf{m}|} \log{\hat{p}_m(d(x, \mathbf{m}_i) \mid l_i)}$\;
LandmarkScore $\gets \sum_{i=0}^{|\mathbf{l}|} \log{\hat{p}_l(d(x, \mathbf{l}_i) \mid \mathbf{v}_i)}$\;
\Return MonitorScore + LandmarkScore
}

\BlankLine
$\mathbf{v} \gets $\{$\mathrm{canberra\_distance}(\mathbf{l}_i, \mathbf{p}^m) \mid \mathbf{l}_i \in \mathbf{l}$\}

$\mathbf{C}$ $\gets$ Constraint-Based-Geolocation($\mathbf{m}$, $\mathbf{p}^m$)\;
$\mathbf{C_l}$ $\gets$ \{$m \in \mathbf{m} \mid \mathbf{C}$ contains $m\} \cup \{l \in \mathbf{l} \mid \mathbf{C}$ contains $l$\}\;
\BlankLine
\Return argmax$_{x\in \mathbf{C_l}}$ Likelihood($x$)

 \caption{Algorithm example}
 \label{algorithm:posit}
\end{algorithm}

\chapter{Background}
\label{sec:background}
\section{Phylogenetic Trees}
    Tree model and how to formulate likelihood on tree space 

\subsection{Evolution Model: Continuous-Time Markov Models}
    What does the (generated) data look like? 
        Generation using evolution models (specifically CTMCs) and transition diagrams 
\subsection{Tree Priors: Yule models}
    Tree priors (specifically Yule and others that I will use) 

\section{Bayesian Phylogenetic Inference}
\label{sec:bayesandmcmc}
\subsection{Likelihood Computation}
\label{sec:likelihoodcomputation}
\subsection{Posterior Sampling with Markov Chain Monte Carlo methods}
\subsection{Implementation in BEAST}

\section{Independent Paradigm}
        likelihood calculations, quality checks and convergence guarantees under such a model 
        
        can make this section + model misspecification its own chapter if its too long
\section{Dependent Paradigm}
\subsection{ZIP Construction}
\subsection{CTMC Product Construction}
Transition diagram
\subsection{Likelihood Computation}

\section{Model Misspecification}
    Model misspecification induced by dependence
        Data generation under P\_dependent but being evaluated on model under P\_independent

\section{Evaluation Metrics}
\subsection{Tree Recovery Metrics}
        Robinson-Foulds (RF) and nRF 
        Clade recovery (over 90\% of trees in posterior) 
        ... and maybe more if time allows 
\subsection{MCMC Diagnostics}
        Autocorrelation, and thus ESS 
        Mixing of chains (visual) 
        Thinning of data perhaps 

\section{ (TBD) helpful things (subequations, algorithms)}
To write equations over multiple lines (like systems of equation or equations too long to fit the page), one can use the \textit{align} environment coupled to the \textit{subequations} environment like so
\begin{subequations}
\begin{align}
    \mathbb{P}(0 \leadsto -a) &= D\int_0^\infty dt\, k_r e^{-k_r t}\int_0^t d\tau\,\partial_x u \big|_{x=-a}, \\
    \mathbb{P}(0 \leadsto b) &= -D\int_0^\infty dt\, k_r e^{-k_r t}\int_0^t d\tau\,\partial_x u \big|_{x=b}.
\end{align}
\end{subequations}

\begin{algorithm}[]
\SetAlgoLined
\SetKwInOut{KwInput}{Input}
\SetKwInOut{KwOutput}{Output}
\SetKwInOut{KwPre}{Pre}
\SetKw{Return}{return}
\SetKwProg{Fn}{Function}{}{end}
\LinesNumbered
\KwInput{$\textbf{m}$, such that $\mathbf{m}_i$ is the position of the $i$'th monitor\newline
$\textbf{l}$, such that $\mathbf{l}_i$ is the position of the $i$'th landmark\newline
$\mathbf{p}^m$, such that $\mathbf{p}^m_i$ is the ping latency from monitor $i$ to the target\newline
$\mathbf{p}^l$, such that $\mathbf{p}^l_i$ is the set of ping latencies to landmark $i$}

\BlankLine
\KwPre{Compute $\hat{p}_m(d \mid l)$, an estimator giving the likelihood of the target being distance $d$ away from the monitor, given that the monitor records a latency of $l$ to that target. Implemented by training a KDE using $\mathbf{p}^l$.\newline
Compute $\hat{p}_l(d \mid v)$, an estimator giving the likelihood of the target being distance $d$ away from the landmark, given a Canberra distance of $v$ between the target and the landmark, using training targets.
}
\BlankLine
\KwOutput{Most likely location of the target}
\BlankLine

\Fn{Likelihood($x$, $\mathbf{v}$)} {
MonitorScore $\gets \sum_{i=0}^{|\mathbf{m}|} \log{\hat{p}_m(d(x, \mathbf{m}_i) \mid l_i)}$\;
LandmarkScore $\gets \sum_{i=0}^{|\mathbf{l}|} \log{\hat{p}_l(d(x, \mathbf{l}_i) \mid \mathbf{v}_i)}$\;
\Return MonitorScore + LandmarkScore
}

\BlankLine
$\mathbf{v} \gets $\{$\mathrm{canberra\_distance}(\mathbf{l}_i, \mathbf{p}^m) \mid \mathbf{l}_i \in \mathbf{l}$\}

$\mathbf{C}$ $\gets$ Constraint-Based-Geolocation($\mathbf{m}$, $\mathbf{p}^m$)\;
$\mathbf{C_l}$ $\gets$ \{$m \in \mathbf{m} \mid \mathbf{C}$ contains $m\} \cup \{l \in \mathbf{l} \mid \mathbf{C}$ contains $l$\}\;
\BlankLine
\Return argmax$_{x\in \mathbf{C_l}}$ Likelihood($x$)

 \caption{Algorithm example}
 \label{algorithm:posit}
\end{algorithm}

\chapter{Methods}
\label{sec:methods}

\section{Simulation Design}
\subsection{Experiment 1} % I should name these things instead of 1 and 2
    Answering (1) % ELABORATE ON DESIGN AS CLEARLY AS POSSIBLE
\subsection{Experiment 2}
    Answering (2) % CAN USE FIGURES TO ILLUSTRATE DESIGN AS CLEARLY AS POSSIBLE
\section{Implementation}
        What data is collected from each run? explain cluster submission 
        How is data compiled and processed? 
\section{Validation of Simulation Framework}
        [1] More (independent) data leads to better recovery (may be worth explaining mathematically why this is true)
        [2] Show that dependent generation is as expected (graphs showing amount of data generated, especially in answering (2))
        [3] Likelihood sanity checks? Probably need to think about this later when I actually do likelihood calculations
        [4] Actually known trees and known data? although this is less useful
\chapter{Results}
\label{sec:results}

Having described the simulation study methods and validated implementation results in Chapter \ref{sec:methods}, we discuss the results obtained from Experiments 1 and 2, in an attempt to answer Questions 1 and 2 as introduced in Section \ref{sec:qns}. This chapter will be split into two sections, presenting the results from each experiment. An in-depth analysis of each experiment's results will be presented in Section \ref{sec:discussion}.

\section{Experiment 1}
    Answering (1) 
        nRF/clade recovery/ESS vs proportion/depth of proportion (tiers) 

\section{Experiment 2}
    Answering (2) 
        nRF/clade recovery/ESS vs tiers 
\chapter{Discussion}
\label{sec:discussion}

\section{Experiment 1 - Impact of Dependence on Reconstruction Accuracy}
\label{sec:discussionexp1}
\section{Experiment 2 - Efficacy of Additional Dependent Data}
\label{sec:discussionexp2}
Answer questions
Discuss tradeoffs 
Discuss interesting results

\section{Impact of Dependence on MCMC Behaviour}
Across both sections

it may not be very important to talk about this. so kiv

\section{Implications for Phylolinguistics}
Implications for phylolinguistics
\chapter{Conclusion}
\label{sec:conclusion}

\section{Summary of Findings}
{\color{red} can include the fact that writers of kra-dai paper were informed and their comments on the figures}
\section{Limitations}
- Not using rate variation
- im sure theres a lot more things i can add once ive fully written this out
\section{Future Work}

\appendix
\chapter{Additional Mathematical Derivations}
dependent likelihood derivations 

ZIP transition probabilities

proofs of properties
\chapter{Additional Mathematical Derivations}

In this appendix, we consider mathematical derivations of results used throughout the thesis. These were derivations that were not relevant to the main thrust of considering the violation of conditional site-independence in models, and thus were omitted from the main text and included here.


\section{Stationary Distribution of CTMC}
\label{app:stationaryprob}
Here we prove that the binary CTMC found in Figure \ref{fig:markovchain} has stationary distribution \begin{equation}
\label{eq:stationarydistributionpi}
\boldsymbol{\pi} = (\pi_0, \pi_1) = \left(\frac{q_{10}}{q_{01}+q_{10}}, \frac{q_{01}}{q_{01}+q_{10}}\right).
\end{equation}

We start by considering $\mathbf{G}$. We note that any stationary distribution $\boldsymbol{\pi}$ satisfies $\boldsymbol{\pi}\mathbf{G} = \boldsymbol{0}$. Clearly, this is true of (\ref{eq:stationarydistributionpi}):
\begin{align}
    \boldsymbol{\pi}\mathbf{G} &= \begin{pmatrix}\frac{q_{10}}{q_{01}+q_{10}} & \frac{q_{01}}{q_{01}+q_{10}}\end{pmatrix} \begin{pmatrix}
        -q_{01} & q_{01} \\ q_{10} & -q_{10}
    \end{pmatrix} \\
    &= \begin{pmatrix}\frac{-q_{10}q_{01} +  q_{01}q_{10}}{q_{01}+q_{10}} & \frac{q_{10}q_{01} - q_{01}q_{10}}{q_{01}q_{10}}\end{pmatrix} = \mathbf{0}.
\end{align}

Thus $\boldsymbol{\pi}$ is a stationary distribution of the CTMC with generator $\mathbf{G}$. We now consider the uniqueness of $\boldsymbol{\pi}$. Clearly both states are connected, making the Markov chain generated from the CTMC irreducible. Since the state space is finite, we have that the CTMC must admit a unique stationary distribution \cite{Douc2019}. This justification is also used for Proposition \ref{prop:ctmcprod} to show that $\tilde{\boldsymbol{\pi}}$ is also stationary.

\section{Markov Property for Stochastic Process $\mathbf{D}_{l, k}$}
\label{app:ctmcprodmarkov}

We now turn to consider the fact that $\mathbf{D}_{l, k}$, the resulting binary encoding of the element-wise product of dependency $\mathbf{D}_{l-1, j}$ and simulated CTMC $\tilde{\mathbf{D}}_{l, k}$ as considered by (\ref{eq:dependencystructure}) in Section \ref{sec:ctmcproduct}, is not in fact a CTMC, because it does not satisfy the Markov property.

We define a partition on state space $\mathcal{S}$:
\begin{equation}
\label{eq:lumpability}
    \mathcal{P} = \{\{00, 01, 10\},\{11\}\}
\end{equation}
This is a clear case of reducing the size of the state space. We employ the idea of \textit{Markov lumpability} \cite{Kemedy1976}:
\begin{definition} [Lumpability]
Suppose $(X_t)_{t\geq 0}$ is a CTMC with finite state space $\mathcal{S}$ and generator matrix $\mathbf{G} = (g_{xy})_{x, y \in \mathcal{S}}$, and $\mathcal{P} := \{P_1, ..., P_m\}$ is a partition of $\mathcal{S}$, then aggregated process $Y_t = i$ where $X_t \in B_i$, then we say that the CTMC is lumpable with respect to partition $\mathcal{P}$ if for every pair of blocks $P_i, P_j$, $\forall x, x' \in P_i$:
\begin{equation}
    \sum_{y\in P_j}g_{xy} = \sum_{y \in P_j}q_{x'y}.
\end{equation}
\end{definition}
This condition can be interpreted as ensuring that the jump probability from block $P_i$ to block $P_j$ is constant, regardless of the state within $P_i$. 

Clearly, this is an important property the Markov chain must have in order for the reduced state space stochastic process to be considered a Markov chain as well. We state the formal result without proof (of which can be found in \cite{Kemedy1976}):
\begin{proposition} [Resulting process of lumped Markov chain is Markovian]
\label{prop:markovprocesses}
Suppose $(X_t)_{t\geq 0}$ is a CTMC with finite state space $\mathcal{S}$ and generator matrix $\mathbf{G} = (g_{xy})_{x, y \in \mathcal{S}}$, and $\mathcal{P} := \{P_1, ..., P_m\}$ is a partition of $\mathcal{S}$. Then the resulting aggregated process $(Y_t)_{t\geq0}$ is also a CTMC if and only if the CTMC is lumpable.
\end{proposition}

As such, in order to prove that $\mathbf{D}_{l, k}$ does not satisfy the Markov property, we need only show that $\mathbf{B}_{l, k}$ is not lumpable with respect to the partition (\ref{eq:lumpability}). This is true, as the generator $\mathbf{G}_{\text{prod}}$ as defined in (\ref{eq:ctmcprodgenerator}) shows: we note that both $\{01, 10\}$ have an equivalent jump probability to $\{11\}$, but that is not the case with state $\{00\}$, which has a zero transition probability to $\{11\}$. As such, from Proposition \ref{prop:markovprocesses}, we can conclude that the lumped chain $\mathbf{D}_{l, k}$ with respect to the partition $\mathcal{P}$ defined in \ref{eq:lumpability} indeed does not follow the Markov property.
\chapter{Simulation Parameters}
\label{app:simulationparameters}

{\color{red} important for sure. need to write it down}

%%%%%%%%%%%%%%%%%%%%%%%%%%%%%%%%%%%%%
\bibliography{report/bibs/bibliography}
%%%%%%%%%%%%%%%%%%%%%%%%%%%%%%%%%%%%%

\end{document}